\section{JUSTIFICACIÓN}

La seguridad pública y la prevención de la violencia representan pilares indispensables para el desarrollo socioeconómico, la gobernabilidad democrática y la calidad de vida en las urbes contemporáneas. La criminalidad impone severos costos directos e indirectos, que comprenden desde la pérdida de vidas y daños materiales hasta la retracción de inversiones, el deterioro del espacio público y la degradación del tejido social.

Contar con un sistema que procese analíticamente cientos de miles de eventos delictivos no estructurados aporta un \textbf{valor práctico directo}: permite a las instituciones policiales y de seguridad ciudadana abandonar el patrullaje reactivo y avanzar hacia un modelo de policía orientada a la solución de problemas (\textit{Problem-Oriented Policing}). Al identificar de manera objetiva los micro-puntos calientes (\textit{hotspots}) y los corredores criminales que cruzan fronteras entre distritos, se optimiza la movilización de personal, se focaliza la instalación de luminarias y cámaras de vigilancia, y se protege con mayor eficacia a las comunidades más vulnerables. Asimismo, en el contexto de Bolivia, donde los datos oficiales de criminalidad no se encuentran centralizados en portales abiertos, la estructuración y análisis de crónicas periodísticas provenientes de los nueve departamentos provee una radiografía inédita y georreferenciada de la conflictividad y el delito en el país.

Desde la \textbf{perspectiva académica y científica}, este proyecto constituye una oportunidad metodológica rigurosa para comparar, en un escenario real de gran escala y alta complejidad dimensional, diversas familias fundamentales de algoritmos de \textbf{Aprendizaje Automático No Supervisado}:
\begin{enumerate}
    \item Algoritmos basados en densidad espacial adaptados a geometrías esféricas terrestres (DBSCAN con Haversine, HDBSCAN jerárquico y OPTICS).
    \item Modelos generativos probabilísticos con asignación difusa (\textit{soft clustering}) mediante Mixturas de Gaussianas (GMM).
    \item Métodos particionales combinados con técnicas de minería de texto y vectorización semántica en alta dimensionalidad (K-Means, TF-IDF, LDA).
    \item Estructuras jerárquicas aglomerativas con minimización de varianza interna (Ward).
    \item Redes neuronales profundas con constricción dimensional (\textit{Autoencoders}) para la detección de eventos anómalos.
    \item Técnicas de reducción de dimensionalidad lineales y no lineales (PCA, TruncatedSVD, PPCA, t-SNE, UMAP).
\end{enumerate}

\subsection{Árbol de Objetivos: Medios y Fines}

El desglose operativo de los medios y fines estratégicos derivados del árbol de objetivos establece la hoja de ruta metodológica del proyecto:

\textbf{Medios (Equivalentes a las causas superadas):}
\begin{itemize}
    \item \textbf{Ingeniería geodésica y circular:} Transformar coordenadas a radianes y aplicar la métrica de \textit{Haversine} para erradicar las distorsiones métricas planas, combinada con la codificación trigonométrica senoidal y cosenoidal de ciclos horarios y semanales.
    \item \textbf{Clustering de densidad arbitraria y multiescala:} Implementar DBSCAN y HDBSCAN para descubrir concentraciones delictivas que siguen arterias viales y separar matemáticamente el ruido estocástico mediante umbrales de densidad y persistencia $\lambda$.
    \item \textbf{Cuantificación continua de riesgo probabilístico:} Ajustar modelos GMM mediante Esperanza-Maximización (EM) para delimitar elipses de covarianza al $95.4\%$ ($2\sigma$) que reflejen la dispersión e incertidumbre espacial de los incidentes.
    \item \textbf{Minería semántica de texto libre:} Construir matrices TF-IDF y modelos generativos de tópicos latentes (LDA) para clasificar descripciones operativas de crímenes sin etiquetas supervisadas previas.
    \item \textbf{Modelado neuronal de anomalías:} Entrenar un Autoencoder no lineal que aprenda la variedad de delitos comunes y alerte sobre incidentes cuyo error de reconstrucción (MSE) supere el percentil crítico.
    \item \textbf{Arquitectura de microservicios y visor 3D inmersivo:} Desplegar una API REST en Flask conectada a un globo terráqueo esférico en WebGL para democratizar el acceso interactivo a los resultados analíticos.
\end{itemize}

\textbf{Fines (Equivalentes a los efectos positivos alcanzados):}
\begin{itemize}
    \item \textbf{Reducción de tiempos de focalización:} Disminuir drásticamente el tiempo requerido por analistas y mandos operativos para detectar focos delictivos emergentes.
    \item \textbf{Optimización presupuestaria y logística:} Incrementar la eficiencia en la asignación de cuadrillas de patrullaje preventivo, minimizando recursos ociosos en zonas de baja incidencia intrínseca.
    \item \textbf{Aprovechamiento integral de fuentes abiertas:} Convertir crónicas de medios bolivianos y repositorios abiertos internacionales en inteligencia procesable.
    \item \textbf{Aislamiento temprano de anomalías delictivas:} Detectar patrones atípicos de alta peligrosidad antes de que se consoliden como dinámicas delictivas crónicas.
    \item \textbf{Interpretabilidad multimodal para tomadores de decisiones:} Proveer tableros y visualizaciones tridimensionales que traduzcan la complejidad algorítmica en conceptos tácticos intuitivos.
\end{itemize}

